\subsection{Voting Theory}

\content{
        \begin{example} Let's look the following vote. \\
        \begin{tabular}{|l|c|c|c|c|}
        \hline
        & 1 & 3 & 3 & 3 \\
        \hline
        1st Choice  & A & A & O & H \\
        \hline
        2nd Choice  & O & H &H&A \\
        \hline
        3rd Choice & H&O&A&O \\
        \hline
        \end{tabular}
        \end{example}
}{% presentation
\vspace{3in}
}{% nopresentation
\vspace{3in}
}{% comments
Do plurality method. 
}

\content{
        \begin{example} 
        Perform the Pairwise comparison method with the following preference
        schedule. \\
        \begin{tabular}{|l|c|c|c|c|} 
        \hline
        & 5&5&6&4\\
        \hline
        1st & D&A&C&B\\
        \hline
        2nd & A&C&B&D\\
        \hline
        3rd & C&B&D&A\\
        \hline
        4th & B&D&A&C\\
        \hline
        \end{tabular}
        \end{example}
}{% presentation
\vspace{4in}
}{% nopresentation
\vspace{4in}
}{% comments
First round: \\
$A: 1.5, B: 1.5, C: 2, D: 1. $\\
If $D$ drops out, that affects who wins!  
Second round: \\
$A: 1.5, B: .5, C: 1. $\\ 
This violates the IIA criterion.
}

\content{
\begin{example} Who wins the following election under the plurality with
elimination method? \\
\begin{tabular}{|l|c|c|c|c|}
\hline
& 37 & 22 & 12 & 29 \\
\hline
1st & A&B&B&C\\
\hline
2nd & B&C&A&A\\
\hline
3rd & C&A&C&B\\
\hline
\end{tabular}
\end{example}
}{% presentation
\vspace{2in}
}{% nopresentation
\vspace{2in}
}{% comments
$C$ is eliminated first, then $A$ wins with the majority. 
}

\content{
Now, suppose that after officials have announced the results, the ballots were
accidentally destroyed before they could be certified, so the votes have to be
recast. Ten of the people whose preference was $B,A,C$ decide that they want to
now support the winner $A$, so they change their votes to $A,B,C$, all the other
votes remain the same. What are the results of the election now? \\
\begin{tabular}{|l|p{0.25in}|p{0.25in}|p{0.25in}|p{0.25in}|}
\hline
&  &  &  &  \\
\hline
1st & &&&\\
\hline
2nd & &&&\\
\hline
3rd & &&&\\
\hline
\end{tabular}
}{% presentation
\vspace{3in}
}{% nopresentation
\vspace{3in}
}{% comments
Now $B$ is eliminated first, and $C$ wins.
Discuss the monotonicity criterion, this is a violation of it. 
}

\content{
        \begin{example} Use Borda Count method to determine the winner of the
        election.\\
        \begin{tabular}{|l|c|c|c|c|}
        \hline
        & 51 & 25 & 10 & 14 \\
        \hline
        1st & S&T&P&O\\
        \hline
        2nd & T&P&T&T\\
        \hline
        3rd & O&O&O&P\\
        \hline
        4th & P&S&S&S\\
        \hline
        \end{tabular}
        \end{example}
}{% presentation
\vspace{3in}
}{% nopresentation
\vspace{3in}
}{% comments
After Borda count: \\
$S: 253, T: 325, O: 228, P: 194. $\\
Notice that $T$ wins but someone else has a majority, review majority criterion. 
}

\subsection{Weighted Voting}

\content{
        \begin{example} Determine all the players who have veto power in the
        following weighted voting system: 
        $$ [34: 12, 11, 10, 7, 3 ]. $$
        \end{example}
}{% presentation
\vspace{3in}
}{% nopresentation
\vspace{3in}
}{% comments
Players 1, 2, and 3 all have veto power. 
}

\content{
        \begin{definition} To calculate the Banzhaf Power Index: 
        \begin{itemize}
        \item List all winning coalitions.
        \item In each winning coalition, determine all players who are critical.
        \item Count up how many times each play is critical. 
        \item Convert these to fractions by dividing by the total times any
        player is critical. 
        \end{itemize}
        \end{definition}
}{% presentation
}{% nopresentation
}{% comments
}

\content{
        \begin{example} Calculate the Banzhaf Power Index for the three players
        in the following weighted voting system 
        $$ [10: 7, 6, 2]. $$
        \end{example}
}{% presentation
\vspace{3in}
}{% nopresentation
\vspace{3in}
}{% comments
Players 1 and 2 both have half the power, Player 3 is a dummy.
}

\content{
        \begin{example} Calculate the Banzhaf Power Index for the following
        weighted voting system
        $$ [ 8: 6,3,2 ]. $$
        \end{example}
}{% presentation
\vspace{3in}
}{% nopresentation
\vspace{3in}
}{% comments
$P_1 = 3/5$, $P_2 = 1/5$ and $P_3 =1/5$.
}

\content{
        \begin{definition} (Shapley-Shubik power index) To compute the
        shapley-shubik power index: 
        \begin{itemize}
        \item List all sequential coalitions
        \item Determine all pivotal players
        \item Count up how many times each player was critical
        \item convert these into fractions by dividing each count by the total
        number of sequential coalitions
        \end{itemize}
        \end{definition}
}{% presentation
}{% nopresentation
}{% comments
}

\content{
        \begin{example} Calculate the Shapley-Shubik Power Index for the following
        weighted voting system
        $$ [ 8: 6,3,2 ]. $$
        \end{example}
}{% presentation
\vspace{3in}
}{% nopresentation
\vspace{3in}
}{% comments
Powers are (in order)\\
$1/2, 1/3, 1/6. $
}

\subsection{Apportionment}

\content{
        \begin{definition} For \textbf{Hamilton's Method}, we perform the
        following steps: 
        \begin{itemize}
        \item Calculate the \textbf{divisor}: (total population divided by total
        number of representatives) 
        \item Calculate the \textbf{quota} for each state: (state's population
        divided by divisor)
        \item Calculate \textbf{lower quotas}
        \item Assign any remaining representatives to those states whose decimal
        parts of the quota was largest. 
        \end{itemize}
        \end{definition}
}{% presentation
}{% nopresentation
}{% comments
lower quotas are rounded down. 
}

\content{
\begin{example} The state of Delaware has three counties: Kent, New Castle, 
and Sussex. The Delaware state House of Representatives has 41 members. If 
Delaware wants to divide this representation along county lines,
let’s use Hamilton’s method to 
apportion them. The populations of the counties are as follows (from the 
2010 Census): Kent: 162,310, New Castle: 538,479, and Sussex: 197,145.
\end{example}
}{% presentation
\vspace{3in}
}{% nopresentation
\vspace{3in}
}{% comments
total pop is 897,934.
}

\content{
        \begin{definition} Jefferson's method:
        \begin{itemize}
        \item Find the Standard divisor
        \item Find each states lower quota
        \item If the lower quotas do not add up to the number of seats
        available, then adjust the divisor, and go back to step 2. 
        \end{itemize}
        \end{definition}
}{% presentation
}{% nopresentation
}{% comments
}

\content{
        \begin{example} Let's redo the first example from last time, this time
        with Jefferson's method. \\
        \begin{tabular}{lr}
        County & Population \\
        \hline
        Kent & 162,310 \\
        New Castle & 538,479 \\
        Sussex & 197,145 \\
        \hline
        Total & 1,052,567
        \end{tabular}
        \end{example}
}{% presentation
\vspace{3in}
}{% nopresentation
\vspace{3in}
}{% comments
}

\subsection{Fair Division}

\content{
        \begin{example} Suppose that Betty and Cade are sharing a cake which is
        half strawberry and half vanilla. If there are six slices of cake total,
        and cade values straberry 3 times as much as he values vanilla, how much
        is each slice of straberry worth to him? 
        \end{example}
}{%presentation
\vspace{3in}
}{% nopresentation
\vspace{3in}
}{% comments
}

\content{
        \begin{example} Troy and Abed are sharing a Subway sandwich, which is
        a foot long and half veggie and half chicken. Troy likes chicken 2 times as much as
        veggies, and Abed likes chicken and veggies equally. If Troy is the
        divider, describe how he will split the sandwich (only vertical slices
        are allowed). Which part will Abed choose? 
        \end{example}
}{%presentation
\vspace{4in}
}{% nopresentation
\vspace{4in}
}{% comments
}

\content{
        \begin{definition} (The Lone Divider Method) 
        \begin{enumerate}
        \item The divider divides the item into $N$ pieces, which we'll label
        $s_1$, $s_2$, ..., $s_N$.
        \item Each of the choosers will separately list which pieces they
        consider to be a fair share. This is called their bid. 
        \item The lists are examined. There are two possibilities: 
        \begin{enumerate}
            \item If it is possible to assign each player a piece that they bid
            on, then do so, and the divider gets the remaining piece. 
            \item If this is not possible, then the divider receives an
            uncontested piece, and the remaining $N-1$ parties repeat the entire
            procedure. If there are only 2 players left, they may use the
            Divider Chooser method. 
            \end{enumerate}
        \end{enumerate}
        \end{definition}
}{%presentation
}{% nopresentation
}{% comments
}

\content{
        \begin{example} Consider the following example. Who is the divider? Is
        there a fair division for this example? Why does the Lone Divider method
        guarantee a fair division in this example? \\
        \begin{tabular}{|l|l|l|l|}
        \hline
        & Piece 1 & Piece 2 & Piece 3 \\
        \hline
        Chad & 40\% & 30 \% & 30 \% \\
        \hline
        Brody & 45\% & 30\%  & 25\% \\
        \hline
        Mary & 33.3\% & 33.3\% & 33.3\% \\
        \hline
        \end{tabular}
        \end{example}
}{%presentation
\vspace{2in}
}{% nopresentation
\vspace{2in}
}{% comments
}

\content{
        \begin{definition} The method begins by compiling a list of items to be
        divided. Then: 
        \begin{enumerate}
        \item Each party involved lists, in secret, a dollar amount they 
        value each item to be worth. This is their sealed bid.
        \item The bids are collected. For each party, the value of all the 
        items is totaled, and divided by the number of parties. This 
        defines their fair share.  
        \item Each item is awarded to the highest bidder.
        \item For each party, the value of all items received is totaled. 
        If the value is more than that party’s fair share, they pay the 
        difference into a holding pile. If the value is less than that party
        ’s fair share, they receive the difference from the holding pile. 
        This ends the initial allocation.  
        \item In most cases, there will be a surplus, or leftover money, in 
        the holding pile.  The surplus is divided evenly between all the 
        players. This produces the final allocation.
        \end{enumerate}
        \end{definition}
}{%presentation

}{% nopresentation
}{% comments
}

\content{
        \begin{example} Jamal, Maggie, and Kendra are dividing an estate 
        consisting of a house, a vacation home, and a small business. Their 
        valuations (in thousands) are shown below. Determine the final 
        allocation.  \\

        \begin{tabular}{|l|l|l|l|}
        \hline
        & Jamal & Maggie & Kendra \\
        \hline
        House & \$250 & \$300 & \$280 \\
        \hline
        Vacation Home & \$170 & \$180 & \$200 \\
        \hline
        Small Business & \$300 & \$255 & \$270 \\
        \hline
        \end{tabular}
        \end{example}
}{%presentation
\vspace{3in}
}{% nopresentation
\vspace{4in}
}{% comments
}
